% Lazo de control dentro del lote con la magnitud que pasa por cada flecha:
% potencia (gruesa), señal eléctrica (fina), magnitud mecánica o material
% (doble) y realimentación (discontinua). Valores típicos, no seleccionados.
\begin{tikzpicture}[
    >=Latex,
    font=\small,
    bloque/.style={rectangle, draw, rounded corners=2pt, minimum height=1.1cm,
        text width=2.3cm, align=center, inner sep=3pt},
    potencia/.style={-{Latex[width=7pt,length=6pt]}, line width=1.6pt},
    senal/.style={-{Latex[width=6pt,length=5pt]}, thick},
    mecanica/.style={-{Latex[width=7pt,length=6pt]}, double, double distance=1.6pt, thick},
    realim/.style={-{Latex[width=6pt,length=5pt]}, thick, dashed},
    nota/.style={font=\footnotesize, align=center}
]

% Cadena principal
\node[nota, text width=1.6cm] (cons) {Consigna\\(pantalla)};
\node[bloque, right=1.6cm of cons, text width=2.7cm] (ctrl) {Controlador\\(dos~velocidades\\y corte)};
\node[bloque, right=2.2cm of ctrl] (drv) {Driver\\del motor};
\node[bloque, right=2.2cm of drv] (mot) {Motor};
\node[bloque, right=1.9cm of mot] (tor) {Tornillo\\y tolva};
\node[bloque, right=1.9cm of tor] (rec) {Recipiente\\sobre celda\\de carga};

\draw[senal] (cons) -- (ctrl)
    node[midway, above, nota] {digital\\0--5 V};
\draw[senal] (ctrl) -- (drv)
    node[midway, above, nota] {PWM o pasos\\0--5 V, mA};
\draw[potencia] (drv) -- (mot)
    node[midway, above, nota] {24 V DC\\1--3 A};
\draw[mecanica] (mot) -- (tor)
    node[midway, above, nota] {par (N\,m)\\y rpm};
\draw[mecanica] (tor) -- (rec)
    node[midway, above, nota] {caudal\\(g/s)};

% Alimentación eléctrica
\node[bloque, above=1.5cm of drv] (fte) {Fuente\\24 V DC};
\node[bloque] (reg) at (ctrl |- fte) {Regulador\\5 V DC};
\node[nota, text width=2.0cm] (red) at (mot |- fte) {Red\\220 V AC};
\draw[potencia] (red) -- (fte);
\draw[potencia] (fte) -- (reg)
    node[midway, above, nota] {24 V DC};
\draw[potencia] (fte) -- (drv)
    node[midway, right, nota] {24 V DC\\hasta 3 A};
\draw[senal] (reg) -- (ctrl)
    node[midway, right, nota] {5 V DC\\$<$0,5 A};

% Medición y realimentación
\node[bloque, below=1.5cm of rec] (adc) {Amplificador\\y ADC\\(HX711)};
\draw[senal] (rec) -- (adc)
    node[midway, right, nota, text width=2.4cm] {mV diferencial\\(0--10 mV)};
\draw[realim] (adc.west) -| (ctrl.south)
    node[pos=0.3, above, nota] {masa medida: digital serie, 24 bits, 10 u 80 lecturas/s};
\node[nota, anchor=east] at ([xshift=-4pt, yshift=-0.75cm]rec.south west)
    {excitación 5 V\\desde el HX711};

% Leyenda
\draw[potencia] ([yshift=-3.9cm]cons.south west) -- ++(0.9,0)
    node[right, nota, align=left] {potencia};
\draw[senal] ([yshift=-3.9cm, xshift=3.0cm]cons.south west) -- ++(0.9,0)
    node[right, nota, align=left] {señal eléctrica};
\draw[mecanica] ([yshift=-3.9cm, xshift=6.6cm]cons.south west) -- ++(0.9,0)
    node[right, nota, align=left] {magnitud mecánica o material};
\draw[realim] ([yshift=-3.9cm, xshift=12.0cm]cons.south west) -- ++(0.9,0)
    node[right, nota, align=left] {realimentación};

\end{tikzpicture}
